\chapter{Threat Model and Evaluation Philosophy}
\label{chap:threat-model}

This chapter makes explicit the threat model used throughout the thesis. Adversarial robustness is meaningful only after the attacker, defender, protected object, perturbation set, and evaluation metric are specified. Ambiguity about any of these elements can make a robustness claim appear stronger than it is. The chapter therefore defines the assumptions behind the reported experiments and explains how those assumptions should be interpreted for foundation vision encoders.

\section{Protected Object}

The protected object is not a single classifier. It is a reusable SSL encoder together with the downstream pipelines that depend on it. This distinction is central. If a model is deployed only as a classifier, then classification robustness is the main property of interest. If a model is deployed as a foundation encoder, then robustness of the shared representation and robustness of downstream task heads both matter.

The encoder can be protected in two senses. First, its representation can be stable: small input perturbations should not cause large changes in encoder space. Second, downstream predictions can be stable: small input perturbations should not cause large changes in task outputs. These two senses are related but not equivalent. The thesis treats them as separate evaluation targets.

\section{Attacker Knowledge}

The experiments use white-box attacks. The attacker knows the encoder, the downstream head when applicable, the loss function, and the input preprocessing. This is a strong but standard assumption in adversarial robustness research. It prevents security through obscurity and tests whether the model itself resists gradient-based perturbations.

White-box evaluation is especially important for foundation models. A foundation encoder may be distributed publicly, made available through libraries, or replicated by downstream users. Even if one deployment hides a specific head, attackers can often approximate it. A defense that works only because the attacker lacks gradients is weaker than a defense that withstands white-box evaluation.

This does not mean black-box attacks are irrelevant. Query-based or transfer attacks are important in deployed systems. However, white-box attacks are a necessary first test. If a model fails severely when the attacker has full access, then claiming robust foundation behavior would be premature.

\section{Attacker Goals}

The thesis considers untargeted degradation rather than targeted manipulation. For classification, the attacker wants the classifier to be wrong. For segmentation, the attacker wants to reduce mIoU by increasing pixel-wise errors. For depth, the attacker wants to increase depth loss and RMSE. For \EmbedAttack, the attacker wants to move the representation away from the clean representation.

Targeted attacks would be a stricter or different setting. A targeted classifier attack might force one specific class. A targeted segmentation attack might turn road pixels into sidewalk pixels. A targeted depth attack might make an obstacle appear farther away. These are important for safety-critical applications, but the supplied resources focus on untargeted robustness. Untargeted attacks are already sufficient to reveal the transfer failure studied here.

The difference between representation and task goals is the core of the thesis. \EmbedAttack has no direct notion of a wrong label, wrong mask, or wrong depth. It only maximizes representation change. Task-specific attacks encode the downstream failure directly. A robust foundation encoder should ideally reduce both kinds of failure.

\section{Perturbation Set}

The experiments use an $\ell_\infty$ perturbation budget. This means every pixel channel can change by at most $\varepsilon$. The default value is $8/255$, a common budget in image robustness work. The budget is small relative to the full pixel range, so perturbations are intended to be visually subtle.

The $\ell_\infty$ set is not the only possible perturbation model. $\ell_2$ perturbations, spatial transformations, patches, corruptions, weather effects, and physical-world changes all matter. However, the $\ell_\infty$ setting is a well-established first test. It also makes the comparison across tasks clean because the same input constraint is used for classification, segmentation, depth, and embedding attacks.

A model that fails at $8/255$ is not necessarily useless, but it is not robust under this standard threat model. A model that succeeds at $8/255$ would still need more tests before deployment. Robustness should be understood as conditional on the perturbation set.

\section{Defender Capabilities}

The defender in this thesis can choose the encoder and robust fine-tuning method. The downstream heads are trained after encoder selection. The evaluated defense, \DeACL, starts from a pre-trained encoder and fine-tunes a robust encoder using representation-level adversarial examples.

The defender does not adversarially train each downstream head against its own task loss in the main evaluation. This is intentional. If every downstream task receives full adversarial training, then the experiment becomes a study of task-specific robust pipelines. The thesis instead asks whether a robust reusable encoder protects multiple downstream tasks.

This defender model is realistic for foundation encoders. A central model provider may release an encoder, and downstream users may train lightweight heads. The provider may not know every future task. Therefore, encoder-level robustness is valuable only if it transfers beyond the exact robustification objective.

\section{Why Task-Specific Adversarial Training Is Not the Baseline}

One might ask why the thesis does not adversarially train segmentation and depth heads directly. The answer is that such training would answer a different question. Task-specific adversarial training can improve task-specific robustness, but it does not prove that the foundation encoder is robust in a task-agnostic way.

The foundation-model promise is amortization. Pre-training costs are paid once, and downstream users benefit many times. If robustness also requires expensive task-specific adversarial training for every downstream task, then robustness has not been amortized. It may still be achievable, but it is no longer a simple property of the shared encoder.

This distinction is important for interpreting \DeACL. \DeACL is attractive precisely because it decouples robustification from downstream labels. The thesis evaluates whether that attractive property is enough. The results suggest that it is not, at least for the studied tasks and attacks.

\section{Evaluation Units}

The evaluation unit is a dataset-task-attack tuple. For example, CIFAR10 under \PGD is one unit, Cityscapes under \SegPGD is another, and NYU-Depth v2 under \DepthPGD is another. This unit matters because results do not necessarily transfer across datasets or attacks.

A robust report should avoid saying simply that an encoder is robust. It should say robust to which attack, on which task, under which metric. The tables in the evaluation chapter are organized exactly this way. They are more verbose than a single score, but the verbosity prevents overclaiming.

This evaluation philosophy also makes failures actionable. If \EmbedAttack performance improves but \SegPGD does not, the next method should incorporate segmentation-aware losses or patch-token constraints. If STL10 improves more than CIFAR100, distribution effects should be investigated. A scalar score would hide these directions.

\section{Clean Utility as a Constraint}

Robustness is not useful if clean performance collapses. The thesis therefore treats clean utility as a constraint, not a secondary detail. Every table reports clean performance next to attacked performance. This reveals trade-offs introduced by robust fine-tuning.

For \DeACL, clean classification drops modestly, clean segmentation drops more visibly, and clean depth worsens. These trade-offs matter. A foundation encoder is shared across many downstream users, so a clean-performance loss affects every user, not only users who require robustness.

A future defense should ideally improve attacked performance while preserving clean utility. If that is impossible, it should at least make the trade-off explicit. Different deployments may choose different points on the trade-off curve, but they need the curve to make an informed decision.

\section{Failure Criteria}

The thesis does not define one universal threshold for failure. Instead, failure is interpreted relative to clean performance and task expectations. If clean CIFAR10 accuracy is 0.91 and attacked accuracy is 0.02, the attack has essentially destroyed the classifier. If clean Cityscapes mIoU is 0.36 and attacked mIoU is 0.03, segmentation is essentially unusable. If clean depth RMSE is 0.68 and attacked RMSE is 1.71, geometry is badly degraded.

This relative interpretation is better than a fixed threshold because tasks use different metrics. Accuracy, mIoU, and RMSE are not directly comparable. What matters is how much the attack changes the task outcome and whether the remaining performance is meaningful.

\section{Adaptive Evaluation}

An evaluation is adaptive when the attack matches the model and task. \EmbedAttack is adaptive to the encoder representation. \PGD is adaptive to the classifier. \SegPGD is adaptive to segmentation. \DepthPGD is adaptive to depth. This is why the thesis reports several attacks instead of one.

A non-adaptive attack can underestimate vulnerability. For example, attacking only the global class token may not fully test a segmentation head that uses patch tokens. Attacking only representation distance may miss a depth-loss direction. The methodology therefore tries to align each attack with the interface being tested.

Adaptive evaluation does not mean unfair evaluation. The perturbation budget remains fixed. The attacker simply optimizes the objective relevant to the deployment. A defense claiming downstream robustness should withstand such attacks.

\section{Scope of Claims}

The thesis claims that, in the supplied experiments, representation-level robust fine-tuning does not reliably transfer to classification, segmentation, and depth task attacks. It does not claim that robust SSL is impossible. It does not claim that \DeACL is ineffective for its intended representation attack. It does not claim that all future methods will fail.

This careful scope is important. The results are negative in a specific sense: the attractive assumption of task-agnostic robustness transfer is unsupported. The results are positive in another sense: \DeACL improves \EmbedAttack robustness, and the evaluation protocol identifies exactly where more work is needed.

\section{Practical Deployment Interpretation}

For deployment, the thesis suggests caution. If a provider releases an encoder advertised as robust because it performs well under a representation attack, downstream users should not assume their task is protected. A segmentation user should test segmentation attacks. A depth user should test depth attacks. A classifier user should test classifier attacks.

The provider can help by releasing a broader robustness report. Such a report should include clean and attacked metrics for representative downstream heads. It should also disclose attack parameters, optimization steps, and whether attacks used the downstream loss. Without this information, robustness claims are hard to interpret.

The same principle applies beyond the tasks studied here. Detection, pose estimation, medical imaging, and video understanding all have their own losses and output structures. A robust foundation encoder should be tested with attacks that match those structures before deployment in high-stakes settings.

\section{Summary}

The threat model clarifies the thesis contribution. The protected object is a reusable SSL encoder and its downstream uses. The attacker is white-box and bounded by an $\ell_\infty$ budget. The attacks differ by objective: representation distance, classification loss, segmentation loss, or depth loss. The defender uses \DeACL representation-level robust fine-tuning. The evaluation reports clean utility and attacked performance separately for each task.

Under this threat model, \DeACL improves representation-attack robustness but does not provide general downstream robustness. This conclusion is meaningful because the threat model is explicit. It also points to future work: robust foundation encoders need objectives and evaluations that cover the downstream interfaces where they are reused.
